\section*{Chapitre \ref{ch:b2:linalg} --- Algèbre linéaire}

\begin{solution}{exo:b2:linalg:1}
Trois formes dans un \omterm{def:b2:linalg:dual}{dual} de dimension $3$~: la liberté suffit. Une
relation $\alpha\varphi_1 + \beta\varphi_2 + \gamma\varphi_3 = 0$
évaluée en $(1,0,0), (0,1,0), (0,0,1)$ donne $\alpha + \gamma = 0$,
$\alpha + \beta = 0$, $\beta + \gamma = 0$, d'où $\alpha = \beta =
\gamma = 0$.

Base préduale $(u_1, u_2, u_3)$~: résoudre $\varphi_i(u_j) =
\delta_{ij}$. En écrivant $u_j = (x, y, z)$~: pour $u_1$~: $x + y = 1$,
$y + z = 0$, $x + z = 0$ donne $u_1 = \bigl(\tfrac12, \tfrac12,
-\tfrac12\bigr)$~; symétriquement $u_2 = \bigl(-\tfrac12, \tfrac12,
\tfrac12\bigr)$, $u_3 = \bigl(\tfrac12, -\tfrac12, \tfrac12\bigr)$.
\end{solution}

\begin{solution}{exo:b2:linalg:2}
Première matrice~: la seule permutation à produit non nul envoie $1 \mapsto
3$, $2 \mapsto 2$, $3 \mapsto 1$ --- la \omterm{def:b2:structures:sn}{transposition} $(1\,3)$,
signature $-1$~: \omterm{def:b2:linalg:det}{déterminant} $-abc$.

Seconde~: une permutation à produit non nul ne peut mélanger les deux
blocs (un coefficient les reliant est $0$), donc elle se scinde en une
permutation de $\{1,2\}$ fois une de $\{3,4\}$, et la signature est le
produit des deux signatures~: la somme se factorise en
\[
(ad - bc)(eh - fg) .
\]
Règle générale suggérée (et vraie, même démonstration)~: le \omterm{def:b2:linalg:det}{déterminant}
d'une matrice diagonale par blocs est le produit des \omterm{def:b2:linalg:det}{déterminants} des
blocs.
\end{solution}

\begin{solution}{exo:b2:linalg:3}
$F$ est défini par les deux équations indépendantes $\varphi_1(x,y,z,t)
= x + y - z - t = 0$ et $\varphi_2 = x - 2y = 0$~: par
le~\cref{thm:b2:linalg:annihilator} lu à l'envers, $F^\circ =
\operatorname{Vect}(\varphi_1, \varphi_2)$ --- elles sont dans $F^\circ$
par construction, elles sont libres (non proportionnelles), et $\dim F^\circ
= 4 - \dim F = 4 - 2 = 2$ puisque $\dim F = 2$ (deux équations
indépendantes dans $\R^4$). Base~: $(\varphi_1, \varphi_2)$~; dimensions~:
$2 + 2 = 4$, comme le théorème l'exige.
\end{solution}

\begin{solution}{exo:b2:linalg:4}
Les $\varphi_i$ sont $n + 1$ formes sur un espace de dimension $n+1$~:
la liberté suffit. Si $\sum_i \lambda_i \varphi_i = 0$, évaluer sur
le polynôme de Lagrange $L_j$ des nœuds~: $\lambda_j = 0$. La
base préduale est $(L_0, \dots, L_n)$, puisque $\varphi_i(L_j) =
L_j(a_i) = \delta_{ij}$.

Pour la forme intégrale avec les nœuds $0, \frac12, 1$ sur $\R_2[X]$~:
$\int_0^1 P = \sum_i c_i P(a_i)$ avec $c_i = \int_0^1 L_i$. Calculons~:
$L_0 = 2(X - \tfrac12)(X - 1)$, $\int_0^1 L_0 = \frac16$~; $L_1 =
-4X(X-1)$, $\int_0^1 L_1 = \frac46$~; $L_2 = 2X(X - \tfrac12)$,
$\int_0^1 L_2 = \frac16$. Donc
\[
\int_0^1 P = \frac{1}{6}\Bigl(P(0) + 4P\bigl(\tfrac12\bigr) +
P(1)\Bigr)
\quad (P \in \R_2[X]) :
\]
la règle de Simpson, exacte sur les polynômes de degré $\leq 2$ --- un
énoncé sur les \omterm{def:b2:linalg:dual}{bases duales}.
\end{solution}

\begin{solution}{exo:b2:linalg:5}
$\operatorname{im} u = Ka$ pour un certain $a \neq 0$~; alors $u(x) =
\varphi(x)\,a$ où $\varphi(x)$ est la coordonnée de $u(x)$ sur $a$
--- linéaire en $x$. Trace~: compléter $a = e_1$ en une base~; la
matrice de $u$ a pour colonnes $\varphi(e_j)\,e_1$, donc son seul
coefficient diagonal est $\varphi(e_1) = \varphi(a)$~: $\operatorname{tr} u =
\varphi(a)$. Puis
\[
u^2(x) = \varphi(x)\, u(a) = \varphi(x)\varphi(a)\, a
= (\operatorname{tr} u)\, u(x).
\]
\omterm{def:b2:linalg:det}{Déterminant}, en deux cas. \emph{Si $\varphi(a) \neq 0$~:} prendre une
base quelconque de l'hyperplan $\ker\varphi$ et ajouter $a$. Alors $u$
annule $\ker\varphi$ (là $u(x) = \varphi(x)a = 0$) et $u(a) =
\varphi(a)\,a$~: la matrice de $I + u$ est diagonale, $(1, \dots, 1,\,
1 + \varphi(a))$, donc $\det(I + u) = 1 + \varphi(a) = 1 +
\operatorname{tr} u$. \emph{Si $\varphi(a) = 0$~:} alors $a \in
\ker\varphi$~; prendre une base de $\ker\varphi$ dont le premier vecteur
est $a$, et ajouter un vecteur $b$ avec $\varphi(b) = 1$. Alors $I + u$
fixe la base de $\ker\varphi$ et envoie $b \mapsto b + a$~:
triangulaire de diagonale unité, $\det(I + u) = 1 = 1 +
\operatorname{tr} u$. Les deux cas s'accordent avec la formule.
\end{solution}

\begin{solution}{exo:b2:linalg:6}
Par l'\cref{exo:b2:linalg:9}, l'hyperplan est $H_A = \{M :
\operatorname{tr}(AM) = 0\}$ avec $A \neq 0$.

\emph{Si $A = \lambda I$~:} $H_A$ est l'hyperplan de trace nulle~; la
matrice de la permutation en $n$-cycle (des uns en positions $(i, i+1)$
et $(n, 1)$) est inversible (son \omterm{def:b2:linalg:det}{déterminant} est $\pm 1$ par
le calcul de l'\cref{ex:b2:linalg:permexample}) et est de trace nulle.

\emph{Si $A$ n'est pas scalaire~:} trouver d'abord un $P$ inversible tel que
$B = P^{-1}AP$ ait un coefficient hors diagonale non nul $b_{ji}$ ($j \neq
i$). En effet, si $A$ en a déjà un, prendre $P = I$~; si $A$ est
diagonale avec deux coefficients distincts $d_1 \neq d_2$, conjuguer par
la transvection $P = I + E_{12}$ produit le coefficient hors diagonale
$d_1 - d_2 \neq 0$ (calcul~: $P^{-1}AP = A + (d_1 - d_2)E_{12}$)~;
et une matrice diagonale dont tous les coefficients sont égaux est
scalaire, exclue. Poser maintenant $M' = I + tE_{ij}$ avec $t =
-\operatorname{tr}(B)/b_{ji}$~: alors
\[
\operatorname{tr}(BM') = \operatorname{tr} B + t\,b_{ji} = 0,
\]
et $M'$ est inversible (triangulaire de diagonale unité). En défaisant
la conjugaison, $M = PM'P^{-1}$ est inversible et
$\operatorname{tr}(AM) = \operatorname{tr}(BM') = 0$~: $M \in H_A$.
\end{solution}

\begin{solution}{exo:b2:linalg:7}
$\det(I + tA)$ est, par la formule des permutations, un polynôme en $t$~;
son terme constant est $1$ ($t = 0$). Son coefficient en $t$~: développer
$\det$ comme une \omterm{def:b2:linalg:alternating}{forme alternée} des colonnes $e_j + t\,c_j(A)$~; par
multilinéarité, les termes linéaires en $t$ remplacent exactement un $e_j$ par
$c_j(A)$~:
\[
\sum_{j} \det(e_1, \dots, c_j(A), \dots, e_n)
= \sum_j a_{jj} = \operatorname{tr} A ,
\]
(le \omterm{def:b2:linalg:det}{déterminant} avec toutes les colonnes canoniques sauf $c_j(A)$ à la
place $j$ ramasse le $j$-ème coefficient diagonal). Donc la dérivée en $0$
est $\operatorname{tr} A$.

Soit $g(t) = \det(\eu^{tA})$. La propriété de groupe donne $g(s + t) =
g(s)g(t)$ (multiplicativité de $\det$), $g$ est dérivable, et
$g'(0) = \operatorname{tr} A$ par ce qui précède ($\eu^{tA} = I + tA +
O(t^2)$). Un morphisme dérivable $(\R, +) \to (\R^*, \times)$
vérifie $g' = g'(0)\,g$ (dériver $g(s+t)$ en $s$ en $0$), donc
$g(t) = \eu^{t\operatorname{tr} A}$ par l'unicité des solutions
de $y' = cy$ avec $y(0) = 1$ (volume de première année).
\end{solution}

\begin{solution}{exo:b2:linalg:8}
Soit $v_k = (1, j^k, j^{2k})^{\mathsf T}$ pour $k = 0, 1, 2$. En utilisant $1
+ j + j^2 = 0$ et $j^3 = 1$~:
\[
C v_k =
\begin{pmatrix}
a + b j^k + c j^{2k}\\
c + a j^k + b j^{2k}\\
b + c j^k + a j^{2k}
\end{pmatrix}
= (a + b j^k + c j^{2k})
\begin{pmatrix} 1\\ j^k\\ j^{2k}\end{pmatrix},
\]
(vérifier la deuxième ligne~: $j^k(a + bj^k + cj^{2k}) = aj^k + bj^{2k} +
cj^{3k} = c + aj^k + bj^{2k}$). Donc $v_k$ est un vecteur propre de
valeur propre $\lambda_k = a + bj^k + cj^{2k}$. Les $v_k$ forment une base
(Vandermonde des $1, j, j^2$ distincts), donc $C$ est diagonalisable
avec ces valeurs propres et
\[
\det C = \lambda_0\lambda_1\lambda_2
= (a+b+c)(a + bj + cj^2)(a + bj^2 + cj).
\]
\end{solution}

\begin{solution}{exo:b2:linalg:9}
L'application $\Theta \colon A \mapsto \operatorname{tr}(A\,\cdot)$ est
linéaire de $\mathcal{M}_n(K)$ dans son \omterm{def:b2:linalg:dual}{dual}, entre espaces de même
dimension $n^2$~: l'injectivité suffit. Si $\operatorname{tr}(AM) =
0$ pour tout $M$, prendre $M = E_{ji}$~: $\operatorname{tr}(A E_{ji}) =
a_{ij} = 0$ pour tous $i, j$~: $A = 0$. Donc $\Theta$ est un isomorphisme.

Unicité de la trace (\cref{prop:b2:linalg:trace})~: une forme $t$
annulant tous les commutateurs est $\operatorname{tr}(A\,\cdot)$ pour un
certain $A$ avec $\operatorname{tr}(A(MN - NM)) = 0$ pour tous $M, N$,
c'est-à-dire $\operatorname{tr}((AM - MA)N) = 0$ pour tout $N$ (cyclicité),
c'est-à-dire $AM = MA$ pour tout $M$ (injectivité de $\Theta$)~: $A$ commute
avec tout, donc est scalaire ($A$ commute avec tous les $E_{ij}$ force
les coefficients hors diagonale à $0$ et les coefficients diagonaux à
être égaux), donc $t = c \operatorname{tr}$.
\end{solution}

\begin{solution}{exo:b2:linalg:10}
Travailler sur $\C$ (une matrice réelle est nilpotente ssi elle l'est en
tant que matrice complexe~: la nilpotence est $u^n = 0$).

\emph{Étape 1~: $\operatorname{tr}(u^k) = 0$ pour $k \geq 1$.} Par
récurrence, $u^k v - v u^k = k\, u^k$~: pour $k = 1$ c'est
l'hypothèse~; pour le passage,
\[
u^{k+1}v - vu^{k+1} = u^k(uv - vu) + (u^k v - v u^k)u
= u^{k+1} + k\,u^{k+1} .
\]
En prenant les traces~: $0 = \operatorname{tr}(u^k v) -
\operatorname{tr}(vu^k) = k \operatorname{tr}(u^k)$, donc
$\operatorname{tr}(u^k) = 0$.

\emph{Étape 2~: des traces de puissances nulles impliquent la nilpotence
(sur $\C$).} Soient $\lambda_1, \dots, \lambda_r$ les valeurs propres
non nulles distinctes de $u$ de multiplicités $m_1, \dots, m_r$ (dans le
polynôme caractéristique, qui est scindé sur $\C$ ---
\cref{ch:b2:reduction}). Les traces de puissances sont $\operatorname{tr}(u^k) =
\sum_i m_i \lambda_i^k$ (trigonaliser~: la diagonale de la $k$-ème
puissance d'une matrice triangulaire est faite des $k$-èmes puissances).
Le système $\sum_i m_i \lambda_i^k = 0$ pour $k = 1, \dots, r$ est
inversible au sens de Vandermonde dans les inconnues $m_i\lambda_i$
(matrice $(\lambda_i^{k-1})$ fois la diagonale $\lambda_i$, tous les
$\lambda_i \neq 0$ distincts)~: tout $m_i \lambda_i = 0$, impossible
avec $m_i \geq 1$ sauf si $r = 0$. Donc $u$ n'a aucune valeur propre non
nulle~: son polynôme caractéristique est $(-X)^n$, et Cayley--Hamilton
(\cref{ch:b2:reduction}) donne $u^n = 0$~: nilpotent.
\end{solution}

\begin{solution}{exo:b2:linalg:11}
Notons $V(a_0, \dots, a_n)$ le \omterm{def:b2:linalg:det}{déterminant} et raisonnons par récurrence
sur $n$~; $V(a_0) = 1$ amorce. Fixer $a_0, \dots, a_{n-1}$ et considérer
$D(T) = V(a_0, \dots, a_{n-1}, T)$, le \omterm{def:b2:linalg:det}{déterminant} de dernière colonne
$(1, T, \dots, T^n)$~: en développant suivant cette colonne, $D$ est un
polynôme de degré $\leq n$ en $T$ dont le coefficient de $T^n$ est le
mineur $V(a_0, \dots, a_{n-1})$. Supposons d'abord que $a_0, \dots,
a_{n-1}$ sont distincts. Pour chaque $T = a_i$ ($i < n$) deux colonnes
coïncident, donc $D(a_i) = 0$~: avec $n$ racines distinctes et degré
$\leq n$,
\[
D(T) = V(a_0, \dots, a_{n-1}) \prod_{i=0}^{n-1}(T - a_i),
\]
et $T = a_n$ plus l'hypothèse de récurrence donnent la formule du
produit. Si deux des $a_0, \dots, a_{n-1}$ coïncident, les deux membres
sont $0$ (colonnes répétées~; facteur répété), et la formule est vraie
trivialement.
\end{solution}

\begin{solution}{exo:b2:linalg:12}
\emph{$D$ inversible.} Multiplier à droite par la matrice par blocs
$T = \begin{psmallmatrix} I & 0\\ -D^{-1}C & I\end{psmallmatrix}$,
qui est triangulaire par blocs de diagonale unité, $\det T = 1$ (son
\omterm{def:b2:linalg:det}{déterminant}, par la formule des permutations, ne ramasse que les blocs
diagonaux --- la règle par blocs de l'\cref{exo:b2:linalg:2})~:
\[
\begin{pmatrix} A & B\\ C & D\end{pmatrix} T
= \begin{pmatrix} A - BD^{-1}C & B\\ C - DD^{-1}C & D\end{pmatrix}
= \begin{pmatrix} A - BD^{-1}C & B\\ 0 & D\end{pmatrix},
\]
dont le \omterm{def:b2:linalg:det}{déterminant} est $\det(A - BD^{-1}C)\det D = \det\bigl((A -
BD^{-1}C)D\bigr) = \det(AD - BD^{-1}CD)$. Puisque $CD = DC$,
$BD^{-1}CD = BC$~: le \omterm{def:b2:linalg:det}{déterminant} est $\det(AD - BC)$.

\emph{$D$ général.} Soit $D_t = D + tI$~; alors $CD_t = D_tC$ encore.
Les deux fonctions
\[
f(t) = \det\begin{pmatrix} A & B\\ C & D_t\end{pmatrix}
\qquad\text{et}\qquad
g(t) = \det(AD_t - BC)
\]
sont des fonctions polynomiales de $t$. Le polynôme $\det(D + tI)$ est
unitaire de degré $n$, donc a au plus $n$ racines~: pour tous les $t$
sauf un nombre fini, $D_t$ est inversible et $f(t) = g(t)$ par le
premier cas. Deux polynômes sur un corps infini coïncidant en une
infinité de points sont égaux~: $f = g$, et $t = 0$ conclut.
\end{solution}

\begin{solution}{pb:b2:linalg:1}
\textbf{1.} $\Phi$ est linéaire avec $\ker\Phi = \bigcap_i
\ker\varphi_i$ (un $p$-uplet s'annule ssi chaque entrée s'annule). Pour
les formes coordonnées $\varepsilon_i$ de $K^p$~: $\Phi^{\mathsf
T}(\varepsilon_i) = \varepsilon_i \circ \Phi = \varphi_i$, donc
$\operatorname{im}\Phi^{\mathsf T} \supseteq
\operatorname{Vect}(\varphi_i)$~; réciproquement
$\operatorname{im}\Phi^{\mathsf T}$ est engendré par les
$\Phi^{\mathsf T}(\varepsilon_i)$ (les $\varepsilon_i$ engendrent
$(K^p)^*$). Donc $\operatorname{rk}\Phi = \operatorname{rk}
\Phi^{\mathsf T} = \dim\operatorname{Vect}(\varphi_1, \dots,
\varphi_p) =: r$ (\cref{prop:b2:linalg:transposerank}), et
le théorème du rang donne $\dim\bigcap_i\ker\varphi_i = n - r$.

\textbf{2.} ($\Leftarrow$) Garder une sous-famille libre maximale, disons
$\varphi_1, \dots, \varphi_r$, engendrant le même espace (de sorte que
l'hypothèse se lit encore $\bigcap_{i \leq r}\ker\varphi_i \subseteq
\ker\varphi$~: l'intersection sur tous les $i$ égale celle sur
$i \leq r$, chaque forme écartée étant une combinaison). L'application
$\Psi = (\varphi_1, \dots, \varphi_r) \colon E \to K^r$ est
surjective (question 1~: son rang est $r$). Si $\Psi(x) = \Psi(y)$
alors $x - y \in \ker\Psi \subseteq \ker\varphi$, donc $\varphi(x) =
\varphi(y)$~: $\varphi$ se factorise en $\varphi = \lambda \circ \Psi$
avec $\lambda \colon K^r \to K$ bien défini~; $\lambda$ est linéaire
parce que $\Psi$ est linéaire et surjective (pour $t = \Psi(x)$, $t' =
\Psi(x')$~: $\lambda(t + \alpha t') = \varphi(x + \alpha x') =
\lambda(t) + \alpha\lambda(t')$). En écrivant $\lambda = \sum c_i
\varepsilon_i$~: $\varphi = \sum_{i \leq r} c_i\varphi_i$.
($\Rightarrow$) Si $\varphi = \sum c_i \varphi_i$, tout $x$
annulant chaque $\varphi_i$ annule $\varphi$.

\textbf{3.} Par la question 1, $\dim\bigcap\ker\varphi_i = n - r$
avec $r = \dim\operatorname{Vect}(\varphi_i) \leq p$, et $r = p$
ssi la famille est libre. Un sous-espace $F$ de codimension $p$~: son
\omterm{def:b2:linalg:annihilator}{annulateur} est de dimension $p$
(\cref{thm:b2:linalg:annihilator})~; une base $(\varphi_1, \dots,
\varphi_p)$ de $F^\circ$ donne $F = \bigcap_i\ker\varphi_i$ (la
formule de reconstitution). Moins~: une intersection de $q$ hyperplans a
une dimension $\geq n - q > n - p$ par la question 1.

\textbf{4.} Calculons $\ker\varphi_1 \cap \ker\varphi_2$~: de $x
+ y - z = 0$ et $y + z - t = 0$, paramétrer par $(y, z)$~: $x = z
- y$, $t = y + z$, donnant le plan des vecteurs $(z - y,\; y,\;
z,\; y + z)$. Sur lui, $\psi = x + 2y - t = (z - y) + 2y - (y + z)
= 0$~: par le lemme de factorisation $\psi \in
\operatorname{Vect}(\varphi_1, \varphi_2)$ --- en effet $\psi =
\varphi_1 + \varphi_2$. Mais $\psi' = x + y + t = (z - y) + y +
(y + z) = y + 2z$ n'y est pas identiquement nul ($y = 1, z = 0$
donne $1$)~: $\psi' \notin \operatorname{Vect}(\varphi_1,
\varphi_2)$.

\textbf{5.} Trois formes sur un espace de dimension $3$~: la liberté
suffit. Si $a\psi_0 + b\psi_1 + c\psi_2 = 0$, tester sur $1, X,
X^2$~: $a + b + c = 0$, $b + \frac c2 = 0$, $b + \frac c3 = 0$~;
soustraire les deux dernières donne $c = 0$, puis $b = 0$, $a = 0$.
Base antéduale~: en écrivant $P = \alpha + \beta X + \gamma X^2$ et en
résolvant $\psi_i(P_j) = \delta_{ij}$ ($P(0) = \alpha$, $P(1) =
\alpha + \beta + \gamma$, $\int_0^1 P = \alpha + \frac\beta2 +
\frac\gamma3$)~:
\[
P_0 = 1 - 4X + 3X^2, \qquad
P_1 = -2X + 3X^2, \qquad
P_2 = 6X - 6X^2 .
\]
(Vérification, p.\,ex.~: $\int_0^1 P_2 = 3 - 2 = 1$, $P_2(0) = P_2(1) = 0$.)
Le problème d'interpolation est résolu par les coordonnées dans la
base antéduale~:
\[
P = 1\cdot P_0 + 2\cdot P_1 + \tfrac32\, P_2 = 1 + X
\]
(coefficient de $X$ $-4 - 4 + 9 = 1$, coefficient de $X^2$ $3 + 6 - 9 =
0$)~; en effet $P(0) = 1$, $P(1) = 2$, $\int_0^1 P = \frac32$.

\textbf{6.} Linéarité~: pour tout $\varphi$, $J(x + \alpha
y)(\varphi) = \varphi(x + \alpha y) = J(x)(\varphi) + \alpha
J(y)(\varphi)$, c'est-à-dire $J(x + \alpha y) = J(x) + \alpha J(y)$.
Injectivité~: si $x \neq 0$, compléter $x = e_1$ en une base~; la forme
coordonnée $e_1^*$ vérifie $J(x)(e_1^*) = 1 \neq 0$. Puisque $\dim E^{**} = \dim E^* = \dim
E$, injective implique bijective.

\textbf{7.} Inclusion~: pour $x \in F$ et $\varphi \in F^\circ$,
$J(x)(\varphi) = \varphi(x) = 0$, donc $J(F) \subseteq
F^{\circ\circ}$. Dimensions
(\cref{thm:b2:linalg:annihilator} deux fois)~:
\[
\dim F^{\circ\circ} = \dim E^* - \dim F^\circ
= n - (n - \dim F) = \dim F = \dim J(F),
\]
$J$ étant injective. Donc $J(F) = F^{\circ\circ}$.

\textbf{8.} Première identité~: $\varphi$ annule $F + G$ ssi elle annule
à la fois $F$ et $G$ (elle annule les sommes ssi elle annule les
morceaux)~: $(F+G)^\circ = F^\circ \cap G^\circ$. Seconde~: l'inclusion
$F^\circ + G^\circ \subseteq (F \cap G)^\circ$ est claire (chaque
terme annule $F \cap G$). Dimensions, en utilisant la première identité
et Grassmann~:
\[
\dim(F^\circ + G^\circ) = \dim F^\circ + \dim G^\circ -
\dim(F^\circ \cap G^\circ)
= (n - \dim F) + (n - \dim G) - \bigl(n - \dim(F +
G)\bigr),
\]
ce qui par Grassmann dans $E$ égale $n - \dim(F \cap G) = \dim(F
\cap G)^\circ$~: égalité.

\textbf{9.} Via le lemme~: $\ker\psi \subseteq \ker\varphi$ avec
$p = 1$ donne $\varphi \in \operatorname{Vect}(\psi)$, et
$\varphi \neq 0$ rend le scalaire non nul. Directement~: choisir $x_0$
avec $\psi(x_0) \neq 0$~; tout $x$ s'écrit $x = \bigl(x -
\frac{\psi(x)}{\psi(x_0)}x_0\bigr) + \frac{\psi(x)}{\psi(x_0)}
x_0$ avec le premier terme dans $\ker\psi = \ker\varphi$~; en appliquant
$\varphi$~: $\varphi(x) = \frac{\varphi(x_0)}{\psi(x_0)}\psi(x)$.

\textbf{10.} Prendre la \omterm{def:b2:linalg:dual}{base duale} $(\varphi_1^*, \dots,
\varphi_n^*)$ de $(\varphi_1, \dots, \varphi_n)$ dans $E^{**}$
(\cref{def:b2:linalg:dual} appliquée à $E^*$) et poser $u_j =
J^{-1}(\varphi_j^*)$~: une base de $E$ ($J$ est un isomorphisme,
question 6), avec $\varphi_i(u_j) = J(u_j)(\varphi_i) =
\varphi_j^*(\varphi_i) = \delta_{ij}$. Unicité~: les conditions
$\varphi_i(u_j) = \delta_{ij}$ déterminent $J(u_j)$ sur la base
$(\varphi_i)$, donc déterminent $u_j$.

\textbf{11.} Linéarité~: $(u + \alpha v)^{\mathsf T}\psi = \psi
\circ (u + \alpha v) = u^{\mathsf T}\psi + \alpha\, v^{\mathsf
T}\psi$. Injectivité~: si $u \neq 0$, choisir $x$ avec $u(x) \neq 0$
et $\psi$ avec $\psi(u(x)) \neq 0$ (l'astuce de la forme coordonnée de
la question 6)~: $u^{\mathsf T}\psi \neq 0$. Les espaces $\mathcal{L}(E,F)$
et $\mathcal{L}(F^*, E^*)$ ont tous deux la dimension $\dim E \dim F$~:
bijective. Si $u$ est inversible, la règle de renversement $(vu)^{\mathsf
T} = u^{\mathsf T}v^{\mathsf T}$ donne $u^{\mathsf
T}(u^{-1})^{\mathsf T} = (u^{-1}u)^{\mathsf T} =
\mathrm{id}_{E^*}$ et $(u^{-1})^{\mathsf T}u^{\mathsf T} =
(uu^{-1})^{\mathsf T} = \mathrm{id}_{F^*}$, donc $(u^{\mathsf
T})^{-1} = (u^{-1})^{\mathsf T}$.

\textbf{12.} Pour $x \in E$ et $\psi \in F^*$~:
\[
\bigl(u^{\mathsf T\mathsf T}(J_E x)\bigr)(\psi)
= (J_E x)\bigl(u^{\mathsf T}\psi\bigr)
= (u^{\mathsf T}\psi)(x)
= \psi\bigl(u(x)\bigr)
= \bigl(J_F(u(x))\bigr)(\psi).
\]
Comme $\psi$ est arbitraire, $u^{\mathsf T\mathsf T} \circ J_E = J_F
\circ u$.

\textbf{13.} Par la~\cref{prop:b2:linalg:transposerank}~: $\ker
u^{\mathsf T} = (\operatorname{im} u)^\circ$, donc $u$ surjective
$\iff \operatorname{im} u = F \iff (\operatorname{im}u)^\circ =
\{0\}$ (\cref{thm:b2:linalg:annihilator}) $\iff u^{\mathsf T}$
injective. Et $\operatorname{im} u^{\mathsf T} = (\ker
u)^\circ$, donc $u$ injective $\iff \ker u = \{0\} \iff (\ker
u)^\circ = E^*$ $\iff u^{\mathsf T}$ surjective.

\textbf{14.} Si $u(F) \subseteq F$ et $\varphi \in F^\circ$~:
$(u^{\mathsf T}\varphi)(x) = \varphi(u(x)) = 0$ pour $x \in F$, donc
$u^{\mathsf T}\varphi \in F^\circ$. Réciproquement, si $u(F)
\not\subseteq F$, choisir $x \in F$ avec $u(x) \notin F$~; par la
formule de reconstitution du~\cref{thm:b2:linalg:annihilator} il existe
$\varphi \in F^\circ$ avec $\varphi(u(x)) \neq 0$~: alors
$(u^{\mathsf T}\varphi)(x) \neq 0$ bien que $x \in F$, donc
$u^{\mathsf T}\varphi \notin F^\circ$~: $F^\circ$ n'est pas stable.

\textbf{15.} $u^{\mathsf T} - \lambda\,\mathrm{id}_{E^*} = (u -
\lambda\,\mathrm{id}_E)^{\mathsf T}$ (la \omterm{def:b2:structures:sn}{transposition} est linéaire et
$\mathrm{id}^{\mathsf T} = \mathrm{id}$), donc son noyau est
$(\operatorname{im}(u - \lambda\,\mathrm{id}))^\circ$
(\cref{prop:b2:linalg:transposerank}), de dimension
\[
n - \operatorname{rk}(u - \lambda\,\mathrm{id})
= \dim\ker(u - \lambda\,\mathrm{id})
\]
par le théorème du rang. En particulier un noyau est non nul ssi
l'autre l'est~: mêmes valeurs propres, mêmes multiplicités géométriques.

\textbf{16.} Inclusion~: si $b = u(x)$ et $u^{\mathsf T}\psi =
0$, alors $\psi(b) = \psi(u(x)) = (u^{\mathsf T}\psi)(x) = 0$~: donc
$\operatorname{im} u \subseteq (\ker u^{\mathsf T})_\circ$.
Dimensions~: pour un sous-espace $S \subseteq F^*$, $S_\circ =
J_F^{-1}(S^\circ)$ (dérouler~: $y \in S_\circ$ ssi tout $\psi \in
S$ annule $y$ ssi $J_F(y) \in S^\circ$), donc $\dim S_\circ = \dim
F - \dim S$. Avec $S = \ker u^{\mathsf T}$~:
\[
\dim(\ker u^{\mathsf T})_\circ
= \dim F - \dim\ker u^{\mathsf T}
= \operatorname{rk} u^{\mathsf T} = \operatorname{rk} u :
\]
égalité des dimensions, donc $\operatorname{im} u = (\ker
u^{\mathsf T})_\circ$. Reformulé~: $b \in \operatorname{im} u$ ssi
$\psi(b) = 0$ pour tout $\psi$ avec $u^{\mathsf T}\psi = 0$ ---
l'alternative de Fredholm.

\textbf{17.} Identifier $(K^m)^*$ avec $K^m$ par $y \mapsto \psi_y$,
$\psi_y(v) = y^{\mathsf T}v$~; alors $(u^{\mathsf T}\psi_y)(x) =
y^{\mathsf T}Ax = (A^{\mathsf T}y)^{\mathsf T}x$, donc $u^{\mathsf
T}\psi_y = \psi_{A^{\mathsf T}y}$~: la \omterm{def:b2:linalg:transpose}{transposée} est la matrice
\omterm{def:b2:linalg:transpose}{transposée}. \emph{Au plus une~:} si $Ax = b$ et
$A^{\mathsf T}y = 0$, alors $y^{\mathsf T}b = y^{\mathsf T}Ax =
(A^{\mathsf T}y)^{\mathsf T}x = 0 \neq 1$. \emph{Au moins une~:}
si (i) échoue, la question 16 fournit $\psi_y$ avec $A^{\mathsf T}y
= 0$ et $y^{\mathsf T}b \neq 0$~; réajuster $y$ pour le rendre égal à $1$.

\textbf{18.} $A = \begin{psmallmatrix} 1 & 1 & 0\\ 0 & 1 & 1\\
1 & 2 & 1\end{psmallmatrix}$ (troisième ligne = première + deuxième, donc $A$
est singulière). Résoudre $A^{\mathsf T}y = 0$~: $y_1 + y_3 = 0$, $y_1
+ y_2 + 2y_3 = 0$, $y_2 + y_3 = 0$ donnent $y_1 = y_2 = -y_3$~: la
droite engendrée par $y = (1, 1, -1)$. Fredholm~: résoluble ssi
$y^{\mathsf T}b = b_1 + b_2 - b_3 = 0$, c'est-à-dire $b_3 = b_1 + b_2$
--- visiblement la bonne condition, puisque la troisième équation est la
somme des deux premières.

\textbf{19.} La matrice de $L$ a $1$ sur la diagonale et
$-\frac12$ en positions $(k, k\pm1)$ (mod $n$)~: symétrique, donc
$L^{\mathsf T} = L$ sous l'identification de la question 17.
\emph{Noyau~:} si $Lx = 0$ alors chaque $x_k = \frac12(x_{k-1} +
x_{k+1})$. Soit $k_0$ maximisant $x_k$~; la moyenne des deux
voisins, tous deux $\leq x_{k_0}$, égale $x_{k_0}$ seulement si les deux
égalent $x_{k_0}$~; en se propageant autour du \omterm{def:b2:structures:sn}{cycle}, $x$ est constant.
Réciproquement les constantes sont annulées. Donc $\ker L^{\mathsf T} = \ker L
= \R(1, \dots, 1)$, et l'alternative de Fredholm se lit~: $Lx = b$
résoluble ssi $(1,\dots,1)^{\mathsf T} b = \sum_k b_k = 0$ ---
la condition de compatibilité discrète~: une «~distribution de chaleur~»
sur un anneau peut être réalisée par un potentiel ssi son flux total
s'annule.

\textbf{20.} Si $A$ est antisymétrique et $S$ symétrique~:
\[
\operatorname{tr}(AS) = \operatorname{tr}\bigl((AS)^{\mathsf
T}\bigr) = \operatorname{tr}(S^{\mathsf T}A^{\mathsf T}) =
-\operatorname{tr}(SA) = -\operatorname{tr}(AS),
\]
donc $2\operatorname{tr}(AS) = 0$ et ($\operatorname{char} K \neq
2$) $\operatorname{tr}(AS) = 0$~: $\mathcal{A}_n \subseteq
\mathcal{S}_n^\circ$ (en identifiant le \omterm{def:b2:linalg:dual}{dual} avec les matrices).
Dimensions~: $\dim\mathcal{S}_n^\circ = n^2 - \frac{n(n+1)}2 =
\frac{n(n-1)}2 = \dim\mathcal{A}_n$~: égalité. En échangeant les rôles
(même calcul), $\mathcal{A}_n^\circ = \mathcal{S}_n$.

\textbf{21.} $\operatorname{tr}(I_nM) = \operatorname{tr} M = 0$
pour $M \in \mathfrak{sl}_n$~: la droite $KI_n$ est dans
l'\omterm{def:b2:linalg:annihilator}{annulateur}, dont la dimension est $n^2 - (n^2 - 1) = 1$~: égalité.
Traduit par l'isomorphisme $A \mapsto
\operatorname{tr}(A\,\cdot)$~: une forme s'annulant sur
$\mathfrak{sl}_n$ est $\operatorname{tr}(\lambda I_n\,\cdot) =
\lambda\operatorname{tr}$.

\textbf{22.} Soit $M \in \mathcal{M}_n(K)$. Le polynôme $t
\mapsto \det(M - tI)$ est non nul de degré $n$, donc a au plus
$n$ racines~; $K$ est de caractéristique $0$, donc infini~: choisir
$\lambda \neq 0$ qui n'est pas une racine. Alors $M = (M - \lambda I) +
\lambda I$ écrit $M$ comme somme de deux matrices inversibles.

\textbf{23.} \emph{Étape 1~:} pour $P$ inversible et $X$ arbitraire,
appliquer l'invariance à $M = XP$~: $t(P(XP)P^{-1}) = t(XP)$,
c'est-à-dire $t(PX) = t(XP)$. \emph{Étape 2~:} fixer $X$~; les deux membres de
$t(BX) = t(XB)$ sont linéaires en $B$ et coïncident sur les $B$ inversibles~;
par la question 22 tout $B$ est une somme de deux inversibles, donc ils
coïncident partout. \emph{Étape 3~:} $t$ annule tout commutateur $XB
- BX$~; les commutateurs engendrent $\mathfrak{sl}_n$ (montré dans la
démonstration de la~\cref{prop:b2:linalg:trace}), donc $t$ s'annule sur
$\mathfrak{sl}_n$ et la question 21 donne $t =
c\operatorname{tr}$. (Réciproquement tout $c\operatorname{tr}$ est
invariant par similitude~: la trace est \emph{l}'invariant linéaire de
similitude.)

\textbf{24.} Si $\operatorname{rk} u = r' \leq r$~: prendre une base
$(f_1, \dots, f_{r'})$ de $\operatorname{im} u$ et écrire $u(x)
= \sum_{i=1}^{r'} \psi_i(x) f_i$~; chaque coordonnée $\psi_i(x)$ de
$u(x)$ est linéaire en $x$ (composition de $u$ avec une forme
coordonnée), donc $u$ est une somme de $r' \leq r$ applications de rang
$\leq1$ (compléter par des zéros). Réciproquement, si $u = \sum_{i=1}^r \psi_i(\cdot)f_i$,
alors $\operatorname{im} u \subseteq \operatorname{Vect}(f_1,
\dots, f_r)$~: $\operatorname{rk} u \leq r$. Sous-additivité~: écrire
$u$ avec $\operatorname{rk} u$ termes et $v$ avec
$\operatorname{rk} v$ termes~; la somme a $\operatorname{rk} u +
\operatorname{rk} v$ termes, donc $\operatorname{rk}(u + v) \leq
\operatorname{rk} u + \operatorname{rk} v$.

\textbf{25.} Le dictionnaire~: un sous-espace $F$ correspond à
$F^\circ$ de dimension complémentaire
(\cref{thm:b2:linalg:annihilator}), et retour par la bidualité
(questions 6--7)~; les sommes s'échangent avec les intersections (question 8)~;
une application $u$ correspond à $u^{\mathsf T}$ avec $\ker u^{\mathsf T}
= (\operatorname{im}u)^\circ$, $\operatorname{im}u^{\mathsf T} =
(\ker u)^\circ$, rangs égaux, injectivité/surjectivité
échangées, sous-espaces stables et valeurs propres appariés
(questions 11--15)~; l'équation $u(x) = b$ est
résoluble ssi $b$ est orthogonal à $\ker u^{\mathsf T}$
(questions 16--19)~; et sur $\mathcal{M}_n$ l'accouplement de trace
réalise tout le dictionnaire concrètement, avec la trace comme
unique invariant linéaire de similitude (questions 20--23) et le rang
comme longueur minimale d'une décomposition en tenseurs
élémentaires (question 24). En dimension infinie les décomptes de
dimension échouent et sont remplacés par des hypothèses de
\emph{fermeture} sur les images et par la complétude --- sur les espaces
de Hilbert ceci devient le théorème de représentation de Riesz et la
théorie de Fredholm des opérateurs compacts, démontrée honnêtement dans
le volume de troisième année.
\end{solution}
